% !TEX root = ../00_Main/Main.tex
\documentclass[../00_Main/Main.tex]{subfiles}
\begin{document}

\label{sec:results}

Los valores a los que se hace referencia, salvo que se especifique otra
cosa, corresponden al desempe\~no normalizado $\bar r$
(Ecuaci\'on~\eqref{eq:normalised-severity-materials}) medio de cada
escenario; \emph{referencia} es la condici\'on de entrenamiento y
\emph{generalizaci\'on}, la media en los $21$ escenarios. Los mapas
estad\'isticos por escenario y las matrices de comparaci\'on entre pares
pueden encontrarse en el Ap\'endice~\ref{ap:stat-maps}
(Figuras~\ref{fig:heatmap-welch}--\ref{fig:ensemble-pairwise}).

\subsection{Mejor modelo individual}\label{sec:res-individual}

El Transformer demuestra ser el mejor modelo individual
(Tabla~\ref{tab:global-mean}): $0{,}927$ en la referencia y $0{,}930$ en
generalizaci\'on, frente a $0{,}871$--$0{,}899$ del resto de los modelos
optimizados, con la menor dispersi\'on entre secuencias y la mayor media
en $16$ de los $22$ escenarios (Figura~\ref{fig:heatmap-global}). Su
ventaja es de tama\~no medio frente a DQN ($d = 0{,}50$), al DQN
recurrente ($d = 0{,}57$) y a Q-Learning ($d = 0{,}79$).

\begin{table}[H]
\centering
\caption[Desempe\~no de los modelos individuales]{Modelos individuales:
media en la referencia, desviaci\'on est\'andar entre secuencias
($\sigma$), media en generalizaci\'on, peor escenario y p\'erdida
respecto de la referencia.}
\label{tab:global-mean}
\small
\setlength{\tabcolsep}{4pt}
\begin{tabular}{@{}lccclc@{}}
\toprule
\textbf{Modelo} & \textbf{Ref.} & $\sigma$ & \textbf{Gen.} & \textbf{Peor escenario} & \textbf{P\'erdida} \\
\midrule
Transformer       & $0{,}927$ & $0{,}045$ & $0{,}930$ & Secuencias largas ($0{,}860$) & $0{,}068$ \\
DQN               & $0{,}894$ & $0{,}055$ & $0{,}899$ & Secuencias largas ($0{,}841$) & $0{,}053$ \\
A2C               & $0{,}887$ & $0{,}063$ & $0{,}896$ & Secuencias largas ($0{,}826$) & $0{,}062$ \\
DQN recurrente    & $0{,}899$ & $0{,}049$ & $0{,}889$ & Extrap.\ de severidad ($0{,}831$) & $0{,}068$ \\
Double Q-Learning & $0{,}896$ & $0{,}048$ & $0{,}877$ & Extrap.\ de severidad ($0{,}785$) & $0{,}111$ \\
Q-Learning        & $0{,}887$ & $0{,}061$ & $0{,}871$ & Extrap.\ de severidad ($0{,}762$) & $0{,}124$ \\
Q-Learning base   & $0{,}871$ & $0{,}074$ & $0{,}851$ & Extrap.\ de severidad ($0{,}627$) & $0{,}244$ \\
\bottomrule
\end{tabular}
\end{table}

\begin{landscape}
\begin{figure}[H]
  \centering
  \includegraphics[width=\linewidth]{ind_01_desempeno_por_escenario.png}
  \caption[Desempe\~no medio por escenario (modelos
  individuales)]{Modelos individuales: desempe\~no medio $\bar r$ por
  modelo y escenario. La primera columna es la referencia. Las cinco columnas a la derecha de la l\'inea vertical son las r\'eplicas fuera de muestra.}
  \label{fig:heatmap-global}
\end{figure}
\end{landscape}

En la Figura~\ref{fig:extra-sev-skew}, con las secuencias ordenadas de
menor a mayor $\bar r$, el Transformer queda por encima o al nivel del
resto en los tres escenarios de severidad y en el de longitudes Poisson,
con la mayor ventaja en los peores casos de la distribuci\'on bimodal;
con secuencias largas lo supera el DQN recurrente, y con alta severidad y
longitud $10$ las curvas casi se superponen. Las
Figuras~\ref{fig:c-base}--\ref{fig:c-trf} (Ap\'endice~\ref{ap:per-model}) muestran, para cada modelo, el
desempe\~no de cada secuencia de la referencia, su distribuci\'on y sus
estad\'isticos por bloque.

\begin{figure}[H]
  \centering
  \includegraphics[width=\linewidth]{ind_05_curvas_por_familia.png}
  \caption[Curvas por secuencia en seis escenarios (modelos
  individuales)]{Desempe\~no por secuencia, ordenado de menor a mayor,
  en seis escenarios de generalizaci\'on. Fila superior: severidad;
  fila inferior: longitud y conjunta. El trazo grueso corresponde al
  Transformer.}
  \label{fig:extra-sev-skew}
\end{figure}

\subsection{Ensambles}\label{sec:res-ensembles}

La Tabla~\ref{tab:ensemble-stress} resume los seis ensambles y la
Tabla~\ref{tab:trf-vs-ens} compara cada uno con el Transformer individual;
como los dos grupos se evaluaron por separado, $d$ y $p$ se calculan con la
media, la desviaci\'on est\'andar y el n\'umero de secuencias de cada
escenario.

S\'olo el ensamble ponderado supera al Transformer: $0{,}939$
frente a $0{,}930$ en generalizaci\'on, mayor media en $18$ de los $21$
escenarios y un peor escenario $0{,}040$ m\'as alto. La diferencia es
significativa pero de tama\~no despreciable ($d = 0{,}19$, apenas por
debajo de un efecto peque\~no). En la referencia reduce en un $45\%$ la
distancia al \'optimo que deja Q-Learning,
$(\bar r - \bar r_{\mathrm{ql}})/(1 - \bar r_{\mathrm{ql}})$, frente al
$36\%$ del Transformer. Frente a los dem\'as ensambles, su ventaja es de tama\~no
medio sobre el consenso ($d = 0{,}58$) y despreciable sobre la compuerta
($d = 0{,}17$); ambas son significativas ($p < 10^{-5}$).

\begin{table}[H]
\centering
\caption[Desempe\~no de los ensambles]{Ensambles: media en la referencia,
media en generalizaci\'on, peor escenario y p\'erdida respecto de la
referencia.}
\label{tab:ensemble-stress}
\small
\setlength{\tabcolsep}{4pt}
\begin{tabular}{@{}lcclc@{}}
\toprule
\textbf{Ensamble} & \textbf{Ref.} & \textbf{Gen.} & \textbf{Peor escenario} & \textbf{P\'erdida} \\
\midrule
Ensamble ponderado        & $0{,}937$ & $0{,}939$ & Secuencias largas ($0{,}900$) & $0{,}037$ \\
Compuerta del Transformer & $0{,}928$ & $0{,}931$ & Secuencias largas ($0{,}861$) & $0{,}067$ \\
Consenso con prior        & $0{,}918$ & $0{,}919$ & Secuencias largas ($0{,}870$) & $0{,}047$ \\
Consenso                  & $0{,}889$ & $0{,}905$ & Secuencias largas ($0{,}826$) & $0{,}063$ \\
Voto suave                & $0{,}914$ & $0{,}902$ & Extrap.\ de severidad ($0{,}860$) & $0{,}054$ \\
Voto por acci\'on         & $0{,}914$ & $0{,}901$ & Extrap.\ de severidad ($0{,}859$) & $0{,}055$ \\
\bottomrule
\end{tabular}
\end{table}

\begin{table}[H]
\centering
\caption[Comparaci\'on de cada ensamble con el Transformer]{Cada ensamble
frente al Transformer individual en generalizaci\'on; los valores
positivos indican ventaja del ensamble.}
\label{tab:trf-vs-ens}
\small
\setlength{\tabcolsep}{4pt}
\begin{tabular}{@{}lccccc@{}}
\toprule
\textbf{Ensamble} & \textbf{Diferencia} & $d$ & \textbf{Efecto} & $p$ & \textbf{Mayor media en} \\
\midrule
Ensamble ponderado        & $+0{,}010$ & $0{,}19$  & despreciable & $< 10^{-5}$  & $18$ de $21$ \\
Compuerta del Transformer & $+0{,}001$ & $0{,}02$  & despreciable & $0{,}68$     & $15$ de $21$ \\
Consenso con prior        & $-0{,}010$ & $-0{,}18$ & despreciable & $< 10^{-5}$  & $5$ de $21$ \\
Consenso                  & $-0{,}025$ & $-0{,}39$ & peque\~no    & $< 10^{-22}$ & $3$ de $21$ \\
Voto suave                & $-0{,}027$ & $-0{,}40$ & peque\~no    & $< 10^{-24}$ & $2$ de $21$ \\
Voto por acci\'on         & $-0{,}028$ & $-0{,}42$ & peque\~no    & $< 10^{-26}$ & $2$ de $21$ \\
\bottomrule
\end{tabular}
\end{table}

La compuerta del Transformer obtiene un desempe\~no equivalente al del
Transformer ($d = 0{,}02$; $p = 0{,}68$): supera su media en $15$
escenarios, pero nunca por m\'as de $0{,}006$. Los otros cuatro ensambles
quedan por debajo del Transformer individual. El voto suave y el voto por
acci\'on, que s\'olo promedian o votan ponderando por la confianza, tienen
las medias m\'as bajas (efecto peque\~no, $d \approx -0{,}4$). Su peor
escenario es la extrapolaci\'on de severidad: el $41\%$ y el $42\%$ de sus
secuencias quedan por debajo de $0{,}8$
(Figura~\ref{fig:ensemble-extrapolation}) y su media queda por debajo de
la de tres de sus cuatro miembros: el Transformer ($0{,}996$), A2C
($0{,}929$) y DQN ($0{,}890$).

El mejor ensamble es el ponderado: tiene la mayor media en la
referencia ($0{,}937$) y en generalizaci\'on ($0{,}939$), el peor escenario
m\'as alto ($0{,}900$, secuencias largas) y la menor p\'erdida respecto de
la referencia ($0{,}037$) (Tabla~\ref{tab:ensemble-stress}). Es adem\'as el
ensamble con mayor media en $19$ de los $22$ escenarios; s\'olo lo superan
el consenso, en \texttt{len\_all\_short} y \texttt{joint\_low\_short}, y
la compuerta del Transformer, en \texttt{len\_poisson}
(Figura~\ref{fig:heatmap-ens}).

\subsubsection{Efecto de las reglas fijas sobre la decisi\'on}\label{sec:res-freq}

Tres ensambles agregan reglas fijas al voto: el ensamble ponderado y el
consenso con \emph{prior} lo mezclan con un \emph{prior} de severidad y
aplican una cota de seguridad, y la compuerta del Transformer decide
cu\'ando ejecutar la acci\'on del Transformer; el voto suave, el voto por
acci\'on y el consenso no tienen reglas fijas. Para saber si la decisi\'on
de esos tres ensambles proviene del voto o de las reglas, se repiti\'o su
evaluaci\'on registrando, en cada decisi\'on, la
acci\'on antes y despu\'es de cada regla y la que prefiere el Transformer
(Ap\'endice~\ref{ap:orchestrator}). La repetici\'on reproduce las medias de
la evaluaci\'on principal con diferencias menores que el $0{,}01\%$ de la
escala de $\bar r$ (Tabla~\ref{tab:ens-freq}).



En el consenso con \emph{prior} decide el voto: el \emph{prior} cambia la
acci\'on votada en menos del $5\%$ de las decisiones, tanto en la
referencia como en generalizaci\'on, y el ensamble queda por debajo del
Transformer. En la compuerta decide el Transformer: su acci\'on se ejecuta
en el $97\%$ y el $95\%$ de las decisiones, lo que explica la equivalencia
de sus desempe\~nos. En el ensamble ponderado deciden el voto y el
\emph{prior}: con $\tau = 15$ las distribuciones de los miembros quedan
casi planas (en la referencia, la acci\'on m\'as probable de cada uno tiene
en promedio probabilidad $0{,}19$ y supera a la segunda por
$0{,}02$--$0{,}03$), y el \emph{prior} cambia la acci\'on votada en el
$41\%$ y el $32\%$ de las decisiones. La cota casi nunca act\'ua. Por lo
tanto, en los dos ensambles que igualan o superan al Transformer la
decisi\'on no proviene s\'olo del voto ponderado por confianza, sino de la
regla.

\begin{table}[H]
\centering
\caption[Frecuencia con que las reglas cambian la decisi\'on]{Porcentaje
de las decisiones con recursos disponibles en que ocurre cada evento.}
\label{tab:ens-freq}
\small
\setlength{\tabcolsep}{4pt}
\begin{tabular}{@{}llcc@{}}
\toprule
\textbf{Ensamble} & \textbf{Evento} & \textbf{Ref.} & \textbf{Gen.} \\
\midrule
Compuerta del Transformer & Sigue al Transformer ($g \ge 0{,}5$) & $96{,}9\%$ & $94{,}9\%$ \\
                          & Acci\'on final distinta de la del Transformer & $0{,}8\%$ & $2{,}7\%$ \\
\midrule
Consenso con prior & El \emph{prior} cambia la acci\'on votada & $4{,}5\%$ & $4{,}6\%$ \\
                   & La cota cambia la acci\'on & $0{,}3\%$ & $1{,}5\%$ \\
                   & Acci\'on final distinta de la del Transformer & $62{,}4\%$ & $58{,}4\%$ \\
\midrule
Ensamble ponderado & El \emph{prior} cambia la acci\'on votada & $40{,}9\%$ & $32{,}1\%$ \\
                   & La cota cambia la acci\'on & $0{,}0\%$ & $0{,}3\%$ \\
                   & Acci\'on final distinta de la del Transformer & $61{,}2\%$ & $47{,}4\%$ \\
\bottomrule
\end{tabular}
\end{table}

\begin{landscape}
\begin{figure}[H]
  \centering
  \includegraphics[width=\linewidth]{ens_01_desempeno_por_escenario.png}
  \caption[Desempe\~no medio por escenario (ensambles)]{Ensambles:
  desempe\~no medio $\bar r$ por ensamble y escenario. La primera columna
  es la referencia. Las cinco columnas a la derecha de la l\'inea vertical son las r\'eplicas fuera de muestra.}
  \label{fig:heatmap-ens}
\end{figure}
\end{landscape}

\subsection{Transformer frente al mejor ensamble}\label{sec:res-trf-vs-best}

El ensamble ponderado supera al Transformer en todas las medidas de la
Tabla~\ref{tab:trf-vs-best}: $+0{,}010$ en la media de generalizaci\'on,
con un efecto despreciable ($d = 0{,}19$; Tabla~\ref{tab:trf-vs-ens}),
$+0{,}040$ en el peor escenario y una p\'erdida m\'axima de poco m\'as de la
mitad. El Transformer tiene mayor media s\'olo en \texttt{len\_all\_short}
($0{,}936$ frente a $0{,}908$), en \texttt{joint\_low\_short} ($0{,}969$
frente a $0{,}942$) y, por menos de $0{,}001$, en \texttt{len\_poisson}.
En la extrapolaci\'on de severidad y en la conjunta, el ensamble ponderado
alcanza el \'optimo en todas las secuencias ($\sigma = 0$;
Figura~\ref{fig:ensemble-extrapolation}), pero todas sus decisiones
coinciden con las del Transformer que lo integra: la diferencia con el
Transformer individual proviene de que el ensamble recorta la severidad a
$9$ antes de pas\'arsela a sus miembros. La Figura~\ref{fig:c-ens} muestra
su desempe\~no por secuencia en la referencia, donde tiene la menor
dispersi\'on de todos los sistemas ($\sigma = 0{,}035$).

\begin{table}[H]
\centering
\caption[Transformer frente al ensamble ponderado]{Transformer individual
frente al ensamble ponderado, el mejor ensamble.}
\label{tab:trf-vs-best}
\small
\begin{tabular}{@{}lcc@{}}
\toprule
\textbf{Medida} & \textbf{Transformer} & \textbf{Ensamble ponderado} \\
\midrule
Media en la referencia ($\sigma$) & $0{,}927$ ($0{,}045$) & $0{,}937$ ($0{,}035$) \\
Media en generalizaci\'on         & $0{,}930$ & $0{,}939$ \\
Peor escenario (secuencias largas) & $0{,}860$ & $0{,}900$ \\
Mayor p\'erdida                    & $0{,}068$ & $0{,}037$ \\
Extrapolaci\'on de severidad       & $0{,}996$ & $1{,}000$ \\
Extrapolaci\'on conjunta           & $0{,}997$ & $1{,}000$ \\
Escenarios con mayor media (de $21$) & $3$ & $18$ \\
\bottomrule
\end{tabular}
\end{table}

\vspace{-\baselineskip} % entre flotantes consecutivos, sin el espacio extra
\begin{figure}[H]
  \centering
  \includegraphics[width=\linewidth]{ens_06_curvas_extrapolacion.png}
  \caption[Curvas de los ensambles en los escenarios fuera de
  rango]{Desempe\~no por secuencia, ordenado de menor a mayor, de los seis
  ensambles en los escenarios fuera de rango; trazo grueso: ensamble
  ponderado; l\'ineas punteadas: media de cada ensamble.}
  \label{fig:ensemble-extrapolation}
\end{figure}

\vspace{-\baselineskip} % entre figuras consecutivas, sin el espacio extra
\begin{figure}[H]
  \centering
  \includegraphics[width=0.87\linewidth]{PES_ENS_results.png}
  \caption[Resultados por secuencia: ensamble ponderado]{Resultados por secuencia del ensamble ponderado (\texttt{pes\_ens}) en la referencia.}
  \label{fig:c-ens}
\end{figure}

\subsection{R\'eplicas fuera de muestra}\label{sec:res-heldout}

En las cinco r\'eplicas fuera de muestra, $320$ secuencias nuevas de la
distribuci\'on de la referencia, ning\'un modelo pierde desempe\~no de
forma apreciable (Tabla~\ref{tab:heldout}). La ca\'ida respecto de la
referencia va de $-0{,}004$ a $+0{,}011$, siempre con efecto despreciable
($|d| \le 0{,}15$) y sin significaci\'on ($p \ge 0{,}29$). En los once
modelos optimizados con Optuna queda entre $-0{,}004$ y $+0{,}005$; la
mayor corresponde a Q-Learning base, que no se optimiz\'o, de modo que
diferencias de ese tama\~no son compatibles con la variaci\'on entre
muestras. El ordenamiento se conserva (correlaci\'on de Spearman entre
las medias, $r_s = 0{,}97$): s\'olo cambian de posici\'on modelos cuyas
medias difieren en menos de $0{,}005$ en ambas condiciones. En las mismas $320$ secuencias, el ensamble ponderado
supera al Transformer por $0{,}010$ (prueba $t$ pareada, $p < 10^{-7}$) y
tiene mayor media en las cinco r\'eplicas. Elegir la configuraci\'on
sobre las $64$ secuencias de la referencia no infl\'o, por lo tanto, el
desempe\~no medido en esa condici\'on.

\begin{table}[H]
\centering
\caption[Referencia frente a las r\'eplicas fuera de muestra]{Media en la
referencia y en las cinco r\'eplicas fuera de muestra (entre par\'entesis,
desviaci\'on est\'andar de las medias de las r\'eplicas); ca\'ida:
referencia menos r\'eplicas; $d$: r\'eplicas frente a referencia.}
\label{tab:heldout}
\small
\begin{tabular}{@{}lcccc@{}}
\toprule
\textbf{Modelo} & \textbf{Ref.} & \textbf{Fuera de muestra} & \textbf{Ca\'ida} & $d$ \\
\midrule
Ensamble ponderado          & $0{,}937$ & $0{,}939$ ($0{,}003$) & $-0{,}002$ & $+0{,}04$ \\
Compuerta del Transformer   & $0{,}928$ & $0{,}928$ ($0{,}004$) & $0{,}000$ & $-0{,}01$ \\
Transformer                 & $0{,}927$ & $0{,}929$ ($0{,}003$) & $-0{,}001$ & $+0{,}03$ \\
Consenso con prior          & $0{,}918$ & $0{,}915$ ($0{,}007$) & $+0{,}003$ & $-0{,}07$ \\
Voto por acci\'on           & $0{,}914$ & $0{,}916$ ($0{,}009$) & $-0{,}002$ & $+0{,}04$ \\
Voto suave                  & $0{,}914$ & $0{,}917$ ($0{,}008$) & $-0{,}003$ & $+0{,}06$ \\
DQN recurrente              & $0{,}899$ & $0{,}903$ ($0{,}008$) & $-0{,}004$ & $+0{,}08$ \\
Double Q-Learning           & $0{,}896$ & $0{,}900$ ($0{,}006$) & $-0{,}004$ & $+0{,}07$ \\
DQN                         & $0{,}894$ & $0{,}891$ ($0{,}006$) & $+0{,}002$ & $-0{,}04$ \\
Consenso                    & $0{,}889$ & $0{,}893$ ($0{,}007$) & $-0{,}004$ & $+0{,}06$ \\
A2C                         & $0{,}887$ & $0{,}887$ ($0{,}006$) & $0{,}000$ & $-0{,}01$ \\
Q-Learning                  & $0{,}887$ & $0{,}881$ ($0{,}008$) & $+0{,}005$ & $-0{,}08$ \\
Q-Learning base             & $0{,}871$ & $0{,}860$ ($0{,}010$) & $+0{,}011$ & $-0{,}15$ \\
\bottomrule
\end{tabular}
\end{table}

\subsection{Justificaci\'on a posteriori del ensamble ponderado}\label{sec:res-posthoc}

Como los par\'ametros del ensamble ponderado se fijaron a mano
(Secci\'on~\ref{sec:ens-methods}), se analizaron despu\'es de la
evaluaci\'on, sin modificarlos, de dos maneras: variando cada uno por
separado y contrastando cada regla fija con la asignaci\'on \'optima. Ambos
an\'alisis usan la referencia y las cinco r\'eplicas fuera de muestra
(Ap\'endice~\ref{ap:orchestrator}).

\textbf{Sensibilidad.} La Tabla~\ref{tab:sensitivity} var\'ia un
par\'ametro por vez y deja los dem\'as en sus valores fijos. La
temperatura, el peso y el ancho del \emph{prior} y el peso base del
Transformer tienen su m\'aximo en el valor elegido, tanto en la referencia
como en las r\'eplicas, de modo que la elecci\'on no depende de las $64$
secuencias sobre las que se hizo; pero son m\'aximos angostos: un paso de
la grilla cuesta entre $0{,}002$ y $0{,}008$. La atenuaci\'on $\eta$, el
t\'ermino $0{,}1$ y la cota de seguridad est\'an en una meseta: cambian la
media en a lo sumo $0{,}001$. En las r\'eplicas, sin \emph{prior}
($w_\pi = 0$) el ensamble supera al Transformer por $0{,}004$; con el
\emph{prior} y $\tau = 1$, por $0{,}001$, y con ambos, por $0{,}010$. La
ventaja proviene, por lo tanto, de combinar la temperatura alta con el
\emph{prior}.

\begin{table}[H]
\centering
\caption[Sensibilidad del ensamble ponderado]{Sensibilidad del ensamble
ponderado a cada par\'ametro, con los dem\'as en sus valores fijos. La
\'ultima columna indica los valores con los que supera al Transformer en
las r\'eplicas fuera de muestra por m\'as de dos errores est\'andar
pareados.}
\label{tab:sensitivity}
\small
\begin{tabular}{@{}llllc@{}}
\toprule
\textbf{Par\'ametro} & \textbf{Fijo} & \textbf{Probados} & \textbf{M\'aximo} & \textbf{Supera al Transformer} \\
\midrule
Temperatura $\tau$            & $15$      & $1$ a $50$ ($8$)  & $15$         & $2$ a $30$ \\
Peso del \emph{prior} $w_\pi$ & $0{,}17$  & $0$ a $0{,}5$ ($7$) & $0{,}17$   & $0$ a $0{,}25$ \\
Ancho del \emph{prior} $\sigma_S$ & $3$   & $1$ a $5$ ($5$)   & $3$          & $3$ a $5$ \\
Peso base del Transformer     & $5$       & $1$ a $10$ ($4$)  & $5$          & $2{,}5$ a $10$ \\
Atenuaci\'on $\eta$           & $0{,}3$   & $0{,}1$ a $1$ ($4$) & meseta     & todos \\
T\'ermino $0{,}1 + c_k$       & $0{,}1$   & $0$ a $1$ ($4$)   & meseta       & todos \\
Cota de seguridad             & activa    & s\'i / no          & sin efecto  & ambos \\
\bottomrule
\end{tabular}
\end{table}

\textbf{Contraste con el \'optimo.} La asignaci\'on \'optima de cada ciudad
se reconstruy\'o con la programaci\'on din\'amica que calcula
$S_{\mathrm{mejor}}$ (Ecuaci\'on~\eqref{eq:dp-optimum}), guardando en cada
ciudad la acci\'on que alcanza el m\'inimo, en las $384$ secuencias de la
referencia y de las r\'eplicas; en el $6\%$ de las ciudades, todas con
$S \le 5$, hay m\'as de una acci\'on \'optima, y los resultados no cambian
al considerarlas todas. Mientras queda presupuesto, el \'optimo nunca asigna $0$ y
asigna en promedio entre $0{,}7$ y $1{,}9$ unidades m\'as que la severidad
de la ciudad; cuando el presupuesto se agota, asigna $0$. El \emph{prior}
acierta el ancho pero no el centro: una gaussiana centrada en $S$ ajustada
a esas asignaciones tiene $\sigma_S = 2{,}7$, cercano al $3$ elegido, pero
si tambi\'en se ajusta el centro, este queda en $S + 2{,}2$, con
$\sigma_S = 1{,}4$, y el ajuste mejora mucho; el ancho del \emph{prior}
compensa en parte su centro desplazado. La cota nunca contradice al
\'optimo, que asigna al menos $\lfloor S/2 \rfloor$ siempre que puede,
pero el ensamble casi nunca queda por debajo de ella: cambia $2$ de
$1\,726$ decisiones. La atenuaci\'on $\eta$ va en el sentido del \'optimo,
con un efecto despreciable. El \emph{prior} es la \'unica regla con efecto
apreciable: cambia el $37\%$ de las decisiones con recursos y mejora
$\bar r$ en $0{,}012$ por secuencia.

\textbf{Una regla fija como pol\'itica.} Si el \'optimo asigna algo m\'as
que la severidad, cabe preguntar cu\'anto rinde hacerlo sin ning\'un
modelo. La Tabla~\ref{tab:fixed-rule} eval\'ua la regla
$a = \min(S + k, R)$, que asigna la severidad de la ciudad m\'as $k$
unidades mientras alcancen los recursos, en las mismas secuencias que los
modelos. Con $k = 2$ iguala al ensamble ponderado y supera al Transformer
en $18$ de los $21$ escenarios de generalizaci\'on y en las cinco
r\'eplicas. Pierde con secuencias cortas, en las que le sobra presupuesto
($0{,}845$ frente a $0{,}936$ del Transformer en \texttt{len\_all\_short}),
y con secuencias largas le conviene un $k$ menor. El valor $k = 2$ se
ley\'o de la asignaci\'on \'optima, que usa informaci\'on completa; es
tambi\'en el mejor en la referencia, la condici\'on sobre la que se
ajustaron los modelos. La versi\'on que no requiere conocer el \'optimo,
asignar la severidad ($k = 0$), queda muy por debajo de todos los modelos
entrenados.

\begin{table}[H]
\centering
\caption[Regla fija frente a los modelos]{Desempe\~no medio de la regla
fija $a = \min(S + k, R)$, sin modelo, frente al Transformer y al ensamble
ponderado en las mismas secuencias.}
\label{tab:fixed-rule}
\small
\begin{tabular}{@{}lccc@{}}
\toprule
\textbf{Pol\'itica} & \textbf{Ref.} & \textbf{Gen.} & \textbf{R\'eplicas} \\
\midrule
Regla, $k = 0$ (asigna $S$) & $0{,}782$ & $0{,}805$ & $0{,}787$ \\
Regla, $k = 1$              & $0{,}900$ & $0{,}907$ & $0{,}906$ \\
Regla, $k = 2$              & $0{,}940$ & $0{,}937$ & $0{,}943$ \\
Regla, $k = 3$              & $0{,}939$ & $0{,}940$ & $0{,}941$ \\
\midrule
Transformer                 & $0{,}927$ & $0{,}930$ & $0{,}929$ \\
Ensamble ponderado          & $0{,}937$ & $0{,}939$ & $0{,}939$ \\
\bottomrule
\end{tabular}
\end{table}

\biblio % Needed for referencing to working when compiling individual subfiles - Do not remove
\end{document}
